\section{Introduction}
\label{sec:intro}

Multimodal large language models (MLLMs) have rapidly advanced visual question answering, captioning, and open-ended multimodal reasoning. Yet even strong models remain vulnerable to hallucination: they may generate fluent but weakly grounded statements, invent objects absent from the image, or confidently choose visually unsupported entities under strong language priors. These failures are particularly damaging when the correct evidence is sparse, local, or easily drowned out by semantically irrelevant visual context.

Recent hallucination-mitigation methods have shown that inference-time intervention is a promising alternative to retraining the host model. Contrastive decoding methods such as VCD~\cite{leng2024vcd}, attention-based control methods such as OPERA~\cite{huang2024opera}, and layer-comparison methods such as DoLa~\cite{chuang2024dola} all demonstrate that decoding is a meaningful intervention point. However, these approaches leave unresolved a different and more demanding design question:

\begin{quote}
\textbf{Can one construct a learned grounding-control mechanism that attaches at decode time and remains reusable across heterogeneous MLLM families without model-specific retraining?}
\end{quote}

This question is fundamentally different from improving a single host model. A module may be lightweight yet still be non-portable if its semantics are inherited from the host architecture. In particular, any learned component that consumes model-specific multimodal hidden states and emits corrections directly in model-specific lexical geometry is coupled to the host model, even when the backbone itself is frozen. Under this criterion, the earlier \texttt{Phi\_calib} route is scientifically useful, but it does not satisfy a strict universality requirement.

The central claim of this paper is therefore a \textbf{universality boundary}: strict hot-swappability is impossible when the learned module lives inside model-native hidden, logit, or attention spaces. To make universality a falsifiable scientific claim rather than a loose engineering intuition, the learned component must instead operate in a standardized external space built from frozen public vision embeddings, frozen public text embeddings, decoded candidate spans, and deterministic attachment logic. This observation leads to \textbf{UniGround}, whose learned core $\Psi_{\mathrm{univ}}$ is external to the host MLLM and whose runtime interface is parameter-free.

UniGround is motivated by a more specific insight than generic decode-time correction. When the host model is about to emit a visually grounded object token, the right question is not whether the whole unfinished prefix is globally similar to the image, but whether the \textbf{candidate itself} can wake up the right local evidence from the image and whether that awakened evidence is strong enough to justify intervention. We call this behavior \textbf{candidate-driven visual resonance}. It reframes hallucination mitigation as selective external evidence activation: the controller should intervene only when a visually groundable candidate is risky, and only using evidence awakened by that candidate rather than by unrelated regions or purely textual continuation bias.

This perspective yields two immediate consequences. First, the learned module must be judged not only by its raw mitigation effect, but also by whether the \textbf{same checkpoint} transfers across different host models without learned adaptation. Second, a purportedly token-local universal method must outperform simpler alternatives such as a static global visual prior; otherwise the additional mechanism is not justified.

The paper makes four contributions:
\begin{enumerate}
  \item It formalizes a \textbf{universality boundary} showing why hidden-state-to-lexical-space calibrators cannot support a strict universal plugin claim.
  \item It proposes \textbf{UniGround}, a universal external sidecar that performs candidate-driven visual resonance and conservative object-level grounding control in a shared semantic space.
  \item It introduces a \textbf{parameter-free tokenizer bridge} that attaches the same learned sidecar to heterogeneous MLLMs during decoding without touching host-internal attention.
  \item It defines an \textbf{evaluation protocol} that separates module-level validation from system-level hot-swap results, while auditing global-prior baselines, structure preservation, portability-cost tradeoffs, and explicit failure boundaries.
\end{enumerate}

The remainder of the paper follows this logic. Section~\ref{sec:related} situates UniGround within the contemporary literature. Section~\ref{sec:positioning} positions UniGround relative to the earlier TLRA line and formalizes the universality boundary. Section~\ref{sec:method} presents the universal controller. Section~\ref{sec:eval} specifies the evaluation protocol. Sections~\ref{sec:tables} and~\ref{sec:discussion} organize tables, figures, and the scope of claims.

\section{Related Work}
\label{sec:related}

\subsection{Multimodal hallucination evaluation and benchmark design}

The modern study of multimodal hallucination rests on two complementary evaluation traditions. The first is the object-hallucination line represented by CHAIR~\cite{rohrbach2018chair}, which established the importance of measuring whether generated objects are actually licensed by the image rather than merely rewarded by language overlap. The second is the probing-based LVLM line represented by POPE~\cite{li2023pope}, which reframed hallucination measurement around controlled yes/no object-existence probes and exposed the sensitivity of open-ended generation to object co-occurrence bias. More recent benchmarks such as AMBER~\cite{wang2023amber} extend this perspective toward multi-dimensional evaluation spanning existence, attribute, and relation hallucinations.

These benchmarks are not interchangeable. CHAIR remains useful when the claim concerns object-level faithfulness in descriptive generation, whereas POPE is better aligned with controlled hallucination probing under a shared prompt template. AMBER broadens the scope further. UniGround is designed with this benchmark ecology in mind: its central claim is not tied to a single metric, but its strongest evidence should come from settings where hallucination can be measured under explicit grounding-sensitive protocols rather than solely through generic answer-quality scores.

\subsection{Decode-time mitigation and control}

The closest neighbors to UniGround are inference-time methods that intervene during generation rather than retraining the host model. VCD~\cite{leng2024vcd} contrasts distributions induced by clean and distorted visual inputs in order to suppress language-prior-driven objects. OPERA~\cite{huang2024opera} monitors generation dynamics through attention and retrospection signals, then penalizes trajectories associated with hallucination. DoLa~\cite{chuang2024dola} originates on the LLM side and shows that layerwise logit contrast can improve factual generation without parameter updates; when transferred into multimodal settings, it serves as a strong decode-time baseline rather than a grounding controller in the present sense.

What unifies these methods is not a shared mechanism, but a shared operational insight: decoding is a meaningful intervention surface. What separates them from UniGround is the scientific object being optimized. Existing decode-time methods typically alter generation through contrastive inputs, attention-level penalties, or host-internal distribution comparisons. UniGround instead asks whether a \textbf{single learned controller} can be reused across hosts while remaining external to model-native hidden geometry.

\subsection{Post-hoc correction, model-updating mitigation, and UniGround's position}

Another important family addresses hallucination by changing the response after generation or by updating the model itself. LURE~\cite{zhou2024lure} performs lightweight post-hoc revision grounded in hallucination statistics. Woodpecker~\cite{yin2023woodpecker} decomposes the response into verifiable claims and corrects hallucinations through a multi-stage validation pipeline. Volcano~\cite{lee2024volcano} uses self-feedback and iterative revision, while HA-DPO~\cite{zhao2023hadpo} tackles hallucination through preference-based model updating. Optional text-side verification lineages include SelfCheckGPT~\cite{manakul2023selfcheckgpt}.

These approaches solve a different problem. Post-hoc correction methods can afford richer verification pipelines because they need not operate inside the host decoder's per-step budget; training-based methods can absorb stronger supervision because they are allowed to change model behavior permanently. UniGround targets \textbf{portable decode-time grounding control}: the same learned sidecar should be reusable across host models, no learned per-model attachment is permitted, and any host dependence must be confined to deterministic interface logic.

\section{Positioning and Design Principle}
\label{sec:positioning}

\subsection{Continuity with earlier TLRA}

The earlier TLRA line established three principles that remain central here. First, hallucination mitigation should act \textbf{during generation}, not only through post-hoc filtering. Second, the relevant unit of control is \textbf{candidate-local}. Third, any practical decode-time intervention must preserve linguistic structure rather than indiscriminately suppress visually weak tokens.

UniGround preserves these principles while changing the location and role of the learned component. It is formulated as a decode-time grounding-control interface rather than as a theory of how evidence should be rerouted inside the host MLLM.

\subsection{Why internal lexical alignment is not universal}

A learned module fails to be strictly universal if either its inputs or outputs are host-specific. Universality is broken when:
\begin{enumerate}
  \item the module consumes hidden states whose semantics depend on the host architecture; or
  \item the module emits vectors interpreted directly in the host model's lexical space.
\end{enumerate}

Let $m$ denote a host MLLM with hidden-state space $\mathcal{H}_m$, vocabulary $\mathcal{V}_m$, and unembedding map $U_m$. Suppose a learned controller $f_\theta$ consumes $h_t \in \mathcal{H}_m$ and emits a correction in the host lexical space $\mathbb{R}^{|\mathcal{V}_m|}$. If there is no model-independent isomorphism that simultaneously aligns $\mathcal{H}_m$, $\mathcal{H}_{m'}$, $\mathcal{V}_m$, and $\mathcal{V}_{m'}$ across two hosts $m$ and $m'$, then the same parameters $\theta$ cannot preserve semantic equivalence across both models.

\texttt{Phi\_calib} instantiates this host-coupled pattern: it reads model-native multimodal hidden states and projects them into the active model's lexical geometry.

\subsection{Universality through standardized observation and action spaces}

To make universality defensible, the learned plugin must be defined through model-agnostic observation and action spaces:
\begin{itemize}
  \item \textbf{Observation space:} pixels, frozen image embeddings, frozen text embeddings, decoded strings, and optional frozen region proposals.
  \item \textbf{Action space:} span-level estimates of visual support, contradiction, and abstention.
\end{itemize}

Under this design, the learned plugin never consumes model-native hidden states, never reads \texttt{lm\_head.weight}, and never relies on a learned model-specific bridge.

\section{Methodology}
\label{sec:method}

\subsection{Problem setup and active decoding frontier}

Let $m$ denote a host MLLM with vocabulary $\mathcal{V}_m$. Given an image $I$ and a text prefix $y_{<t}$, the host model produces next-token logits
\[
L_t^{(m)} \in \mathbb{R}^{|\mathcal{V}_m|}.
\]
Rather than overriding the full vocabulary distribution with an external model, UniGround acts on the host model's \textbf{active semantic frontier}
\[
\mathcal{F}_t^{(m)} = \mathrm{Top}\text{-}M\!\left(L_t^{(m)}\right),
\]
with $M$ fixed across runs for comparability. By acting on $\mathcal{F}_t^{(m)}$, UniGround remains a genuine decode-time controller.

We compare two routes: \textbf{\texttt{TLRA\_internal}} (including \texttt{TLRA\_zero} and \texttt{TLRA\_calib}) and \textbf{\texttt{TLRA\_univ}}, the universal sidecar path instantiated by UniGround.

\subsection{Universal observation space}

UniGround is defined over a model-agnostic observation space that is \textbf{candidate-conditioned}. For image $I$, context sketch $c_t$, candidate span $s$, and region bank $\mathcal{R}(I)=\{r_i\}_{i=1}^N$,
\[
\mathcal{O}_{\mathrm{univ}}(I, c_t, s)
= \{E_v(I), E_t(c_t), E_t(s), E_v(r_1), \dots, E_v(r_N)\},
\]
where $E_v$ is a frozen public vision encoder, $E_t$ is a frozen public text encoder, and $c_t = C(y_{<t})$ is a deterministic short context sketch. We write
\[
z_I = E_v(I), \qquad z_c = E_t(c_t), \qquad z_s = E_t(s), \qquad Z_R = \{z_{R_i}\}_{i=1}^N.
\]

\subsection{Span proposal through a parameter-free tokenizer bridge}

Let $B_m$ denote the tokenizer bridge for host $m$. For each candidate token $c \in \mathcal{F}_t^{(m)}$,
\[
s = B_m(c, y_{<t}),
\]
through deterministic steps: decode token or prefix; canonicalize whitespace and subword markers; expand to a short span when permitted; merge equivalent strings. The span set is
\[
\mathcal{S}_t^{(m)} = \{ B_m(c, y_{<t}) \mid c \in \mathcal{F}_t^{(m)} \}.
\]

\subsection{Candidate-driven visual resonance}

Before scoring, UniGround constructs a \textbf{candidate-conditioned evidence summary}. Given span $s$, use $z_s$ to weight regions:
\[
a_{t,i}(s)
:=
\frac{\exp(\cos(z_s, z_{R_i}) / \tau_r)}
{\sum_{j=1}^{N} \exp(\cos(z_s, z_{R_j}) / \tau_r)},
\qquad
\bar z_R(s) = \sum_{i=1}^{N} a_{t,i}(s)\, z_{R_i}.
\]
Given this, $\Psi_{\mathrm{univ}}$ predicts
\[
\mathbf{q}_t(s) = \Psi_{\mathrm{univ}}(z_I, \bar z_R(s), z_c, z_s)
:= \big[\sigma_{\mathrm{sup}}(s), \sigma_{\mathrm{con}}(s), \sigma_{\mathrm{abs}}(s)\big].
\]
UniGround composes the \textbf{grounding-resonance energy}
\[
\mathcal{R}_t(s)
:= g_t(s)\,\big(1-\sigma_{\mathrm{abs}}(s)\big)\,
\big(\sigma_{\mathrm{sup}}(s) - \lambda_{\mathrm{con}}\sigma_{\mathrm{con}}(s)\big),
\]
where $g_t(s) \in [0,1]$ is a deterministic structural gate.

\subsection{Risk-aware activation}

Define the \textbf{risk-aware activation score}
\[
\rho_t(s)
:=
\pi_t^{\mathrm{host}}(s)\,\bigl(1-q_t^{\mathrm{coarse}}(s)\bigr).
\]
The controller activates when
\[
\max_{s \in \mathcal{S}_{t,\mathrm{obj}}^{(m)}} \rho_t(s) > \tau_{\mathrm{risk}},
\]
where $\mathcal{S}_{t,\mathrm{obj}}^{(m)}$ is the object-candidate subset from the bridge and structural gate.

\subsection{Conservative redistribution and abstention-triggered back-off}

Let $\Gamma_t \in \{0,1\}$ indicate whether the risk gate fires at step $t$, i.e., $\max_{s \in \mathcal{S}_{t,\mathrm{obj}}^{(m)}} \rho_t(s) > \tau_{\mathrm{risk}}$. When $\Gamma_t=1$, UniGround applies bounded bias:
\[
b_t(s) = \alpha \, \Gamma_t \, \tanh\!\big(\mathcal{R}_t(s) / \tau_b\big),
\]
\[
\Delta L_t^{(m)}(v)
:=
\sum_{s \in \mathcal{S}_t^{(m)}} A_m(v,s)\, b_t(s),
\qquad
\widetilde{L}_t^{(m)}(v) = L_t^{(m)}(v) + \Delta L_t^{(m)}(v),
\]
where $A_m(v,s)$ is the deterministic attribution map, $\Gamma_t \in \{0,1\}$ indicates activation, $\alpha$ controls strength, and $\tau_b$ bounds bias.

If $\max_{s \in \mathcal{S}_t^{(m)}} \sigma_{\mathrm{abs}}(s) \ge \tau_{\mathrm{abs}}$, UniGround backs off and leaves host logits unchanged.

\subsection{Why the mechanism remains universal}

$\Psi_{\mathrm{univ}}$ operates in the external encoder space; $B_m$ and $A_m$ are parameter-free; the host contributes only its frontier and final decoding.

\subsection{Controls and methodological scope}

UniGround is distinguished from an \textbf{external global prior} (frozen image-level bias without candidate-local verification) and from \texttt{TLRA\_internal} controls (\texttt{TLRA\_zero}, \texttt{TLRA\_calib}, \texttt{Base + LoRA}). The central assertive claim is:

\begin{quote}
A single learned sidecar, defined in a shared external semantic space and attached through parameter-free model interfaces, can exert useful decode-time grounding control across heterogeneous MLLM families without model-specific retraining.
\end{quote}

\section{Evaluation Protocol}
\label{sec:eval}

\subsection{Universal hot-swap validity (H1)}

The same $\Psi_{\mathrm{univ}}$ checkpoint must transfer across MLLM families without learned adaptation: no per-model retraining; no learned per-model adapter; only parameter-free tokenizer or interface logic may vary.

\subsection{Module-level validation (H2)}

Compare \texttt{Psi\_univ}, \texttt{Hardcoded\_Scorer}, and \texttt{Zero\_Scorer}. Report macro-F1, per-class accuracy, abstention behavior, and checkpoint metadata.

\subsection{System-level hallucination reduction (H3)}

Compare Base, VCD, DoLa, \texttt{TLRA\_internal\_zero}, \texttt{TLRA\_internal\_calib}, and \texttt{TLRA\_univ}. Report hallucination metrics, length preservation, intervention coverage, latency, and memory.

\subsection{Structure preservation (H4)}

Compare Base, \texttt{TLRA\_univ}, \texttt{TLRA\_univ\_no\_gate}, and \texttt{TLRA\_internal\_calib}. Audit function-word corruption, continuation stability, and reasoning retention.

\subsection{Portability versus internal ceiling (H5)}

Report the tradeoff if internal calibration achieves higher per-host scores while the universal route preserves hot-swap validity.

\subsection{External global priors (H6)}

Compare \texttt{External\_Global\_Prior}, \texttt{TLRA\_univ\_global\_only}, and full \texttt{TLRA\_univ}.

\subsection{Required ablations}

\texttt{TLRA\_univ\_no\_gate}, \texttt{TLRA\_univ\_global\_only}, \texttt{TLRA\_univ\_region\_only}, \texttt{TLRA\_univ\_no\_abstain}, \texttt{External\_Global\_Prior}, \texttt{TLRA\_internal\_calib}, and \texttt{Base + LoRA} (internal fairness).

\subsection{Failure-boundary audits (H7)}

Audit Top-$M$ ceiling, prefix ambiguity, span collapse, suffix stability, abstention and abort triggers, and latency decomposition.

\section{Main-Paper Tables and Figures}
\label{sec:tables}

Target layout: approximately eight body pages; one strong hero figure after Methodology or at the start of this section; compact Figures~2--3.

\subsection{Figure placeholders}

Figure~\ref{fig:hero} (placeholder) should depict host frontiers, the parameter-free tokenizer bridge, shared $\Psi_{\mathrm{univ}}$, resonance energy $\mathcal{R}_t(s)$, gates, and frontier redistribution---not full-vocabulary reranking.

\begin{figure*}[t]
  \centering
  \fbox{\begin{minipage}{0.95\textwidth}\centering\vspace{36pt}
  \textit{Placeholder: Figure~1 --- Frontier-to-resonance decode-time control (illustrator brief in source benchmark).}\vspace{36pt}
  \end{minipage}}
  \caption{UniGround reads the host's active semantic frontier, lifts spans through a parameter-free tokenizer bridge into a shared external space, scores candidates with one reusable $\Psi_{\mathrm{univ}}$, and redistributes frontier mass with structural gating and abstention-triggered back-off.}
  \Description{Placeholder block diagram for UniGround pipeline.}
  \label{fig:hero}
\end{figure*}

\begin{figure}[t]
  \centering
  \fbox{\begin{minipage}{0.92\columnwidth}\centering\vspace{28pt}
  \textit{Placeholder: Figure~2 --- Portability vs.\ internal-ceiling tradeoff.}\vspace{28pt}
  \end{minipage}}
  \caption{Internal calibration may achieve a higher host-specific ceiling; UniGround targets a stronger portability claim via a shared external controller.}
  \Description{Placeholder conceptual tradeoff figure.}
  \label{fig:tradeoff}
\end{figure}

\begin{figure}[t]
  \centering
  \fbox{\begin{minipage}{0.92\columnwidth}\centering\vspace{28pt}
  \textit{Placeholder: Figure~3 --- Candidate-aware resonance vs.\ global prior.}\vspace{28pt}
  \end{minipage}}
  \caption{Candidate-aware resonance should discriminate competing frontier spans under shared image evidence beyond a static global prior; otherwise the local-control claim contracts.}
  \Description{Placeholder comparison figure.}
  \label{fig:global}
\end{figure}

\subsection{Tables}

\begin{table*}[t]
  \caption{Universal hot-swap grounding results (placeholders).}
  \label{tab:hotswap}
  \centering
  \small
  \begin{tabular}{@{}lccccccc@{}}
    \toprule
    Method & Shared ckpt. & Learned re-attach & Host A & Host B & Host C & Mean & Status \\
    \midrule
    Base & N/A & N/A & TBF & TBF & TBF & TBF & placeholder \\
    VCD & N/A & N/A & TBF & TBF & TBF & TBF & placeholder \\
    DoLa & N/A & N/A & TBF & TBF & TBF & TBF & placeholder \\
    External\_Global\_Prior & N/A & N/A & TBF & TBF & TBF & TBF & placeholder \\
    TLRA\_internal\_zero & N/A & N/A & TBF & TBF & TBF & TBF & placeholder \\
    TLRA\_internal\_calib & no & yes & TBF & TBF & TBF & TBF & placeholder \\
    TLRA\_univ & yes & no & TBF & TBF & TBF & TBF & placeholder \\
    \bottomrule
  \end{tabular}
\end{table*}

\begin{table*}[t]
  \caption{Internal alignment versus universal control (placeholders).}
  \label{tab:internal}
  \centering
  \footnotesize
  \begin{tabular}{@{}lcccccccc@{}}
    \toprule
    Method & Coupling & Hidden & Lexical & Struct.~prot. & Back-off & Ground. & Capab. & ITL \\
    \midrule
    TLRA\_internal\_zero & internal & yes & yes & vocab-side & N/A & TBF & TBF & TBF \\
    TLRA\_internal\_calib & internal & yes & yes & vocab-side & N/A & TBF & TBF & TBF \\
    TLRA\_univ & univ.\ sidecar & no & no & string-side & yes & TBF & TBF & TBF \\
    TLRA\_univ\_no\_gate & univ.\ sidecar & no & no & none & yes & TBF & TBF & TBF \\
    \bottomrule
  \end{tabular}
\end{table*}

\begin{table}[t]
  \caption{Tokenizer-bridge and prefix-stability audit (placeholders).}
  \label{tab:bridge}
  \centering
  \scriptsize
  \begin{tabular}{@{}lccccccc@{}}
    \toprule
    Method & Tok.\ A & Tok.\ B & Pref.\ amb. & Span coll. & Suf.\ stab. & Abst. & Abort \\
    \midrule
    TLRA\_univ & TBF & TBF & TBF & TBF & TBF & TBF & TBF \\
    TLRA\_univ\_no\_gate & TBF & TBF & TBF & TBF & TBF & TBF & TBF \\
    \bottomrule
  \end{tabular}
\end{table}

\section{Discussion and Scope of Claims}
\label{sec:discussion}

The principal contribution is that decode-time grounding control can be externalized into a shared semantic space and remain useful across heterogeneous MLLMs under an explicit portability contract---not that a universal controller must dominate every model-coupled method on every host.

\subsection{When the universal claim is warranted}

A strong claim requires: (i)~same $\Psi_{\mathrm{univ}}$ checkpoint across hosts; (ii)~no learned per-model adapter; (iii)~grounding gains or justified abstention under cost constraints; (iv)~material advantage over external global-prior baselines when claiming candidate-local control.

\subsection{Portability and ceiling}

Stronger per-host results for \texttt{TLRA\_internal\_calib} reveal the expected portability-ceiling tradeoff rather than refuting the universal formulation.

\subsection{When the claim must contract}

If full UniGround matches \texttt{External\_Global\_Prior} or \texttt{TLRA\_univ\_global\_only}, the candidate-aware story weakens. If learned per-model adaptation is required at attachment time, the strict universality claim fails.

\appendix

\section{Universality Boundary Proof Sketch}
Placeholder: hidden-state inputs; lexical-space outputs; parameter-free tokenizer bridges.

\section{\texttt{Psi\_univ} Architecture and Training Contract}
Placeholder: frozen encoders, dimensionality, scorer architecture, objectives.

\section{Tokenizer Bridge and Prefix Ambiguity}
Placeholder: span rules, canonicalization, ambiguity and failure cases.

\section{Internal Control Path}
Placeholder: \texttt{TLRA\_zero}, \texttt{TLRA\_calib}, \texttt{Phi\_calib} as controls.

\section{Structure-Gate Audit}
Placeholder: protected categories, false positives/negatives, continuation stability.

\section{Cost and Failure Boundaries}
Placeholder: Top-$M$, latency decomposition, abstention distribution, global-prior comparison.
